% subsection: 総和型

  \subsection{総和型}
    \subsubsection{総和と一般項の漸化式型 \, $S_n=pa_n+f(n)$}
      ある漸化式によって定められる数列$a_n$によって,
      \[
        S_n=\sum_{k=1}^{n} a_k
      \]
      として定められる数列$S_n$について考えよう.
      $S_n$と$a_n$の関係を表した
      \[
        S_n=pa_n+f(n)
      \]
      のような漸化式を考える.
      このような漸化式は今まで取り扱った形にはない,未知のものである.
      であれば既知の形にするのが数学である.
      これまでのことをきちんと理解していれば,\,$a_{n+1}$と$a_n$の関係が分かっている時には解くことができる.
      $S_n$をうまく変形して$a_{n}$関連だけの式を産むには,\,$p \neq 1$の時,
      \begin{align*}
        &S_n=pa_n+f(n) \tag*{\(\cdots\) (1)}\\
        &S_{n+1}=pa_{n+1}+f(n+1) \tag*{\(\cdots\) (2)}
      \end{align*}
      $(2)-(1)$をして
      \begin{align*}
        &a_{n+1}=pa_{n+1}-pa_n+f(n+1)-f(n)\\
        &a_{n+1}-pa_{n+1}=-pa_n+f(n+1)-f(n)\\
        &a_{n+1}(1-p)=-pa_n+f(n+1)-f(n)\\
        &a_{n+1}=\frac{-pa_n+f(n+1)-f(n)}{(1-p)}\\
        &a_{n+1}=\frac{p}{p-1}a_n+\frac{1}{1-p}\{f(n+1)-f(n)\}
      \end{align*}
      $\displaystyle g(n)=\frac{1}{1-p}\{f(n+1)-f(n)\}$とおくと,
      \begin{align*}
        &a_{n+1}=\frac{p}{p-1}a_n+g(n)
      \end{align*}
      従って漸化式は階差数列型に帰着する.\,$p=1$のときも同様の流れで求めることができる.

      \noindent \textbf{【例題】}
      \begin{enumerate}[label=(\arabic*),parsep=3pt]
        \item $\displaystyle S_n=2a_n-n$
        \item $\displaystyle S_n=\frac12a_n+n$
      \end{enumerate}

      \noindent\textbf{【方針】}
      \begin{enumerate}[label=(\arabic*),parsep=3pt]
        \item まず$n=1$を代入して$a_1$を求めてから,本文の変形をなぞって$a_{n+1}$と$a_n$だけの漸化式を作りましょう.
        \item 同様です.できた漸化式は特性方程式型に帰着します.
      \end{enumerate}

      \noindent\textbf{【解答】}
      \begin{enumerate}[label=(\arabic*),parsep=3pt]
        \item $\displaystyle a_n=2^n-1$
        \item $\displaystyle a_n=1+(-1)^{n-1}$
      \end{enumerate}

      \noindent\textbf{【解法】}
      \begin{enumerate}[label=(\arabic*),parsep=3pt]
      \item $\displaystyle S_n=2a_n-n$

          $n=1$を代入すると,
          \[
            S_1=a_1=2a_1-1
          \]
          より
          \[
            a_1=1
          \]
          また,
          \begin{align*}
            a_{n+1}
              &=S_{n+1}-S_n\\
              &=\{2a_{n+1}-(n+1)\}-(2a_n-n)\\
              &=2a_{n+1}-2a_n-1.
          \end{align*}
          したがって,
          \[
            a_{n+1}=2a_n+1.
          \]
          ここで $b_n=a_n+1$ とおくと,
          \begin{align*}
            b_{n+1}
              &=a_{n+1}+1\\
              &=2a_n+2\\
              &=2(a_n+1)\\
              &=2b_n.
          \end{align*}
          また,
          \[
            b_1=a_1+1=2
          \]
          であるから,
          \[
            b_n=2^n.
          \]
          よって,
          \[
            a_n=2^n-1.
          \]
      \item $\displaystyle S_n=\frac12a_n+n$
          $n=1$を代入すると,
          \[
            S_1=a_1=\frac12a_1+1
          \]
          より
          \[
            a_1=2
          \]
          \begin{align*}
            a_{n+1}
              &=S_{n+1}-S_n\\
              &=\left\{\frac12a_{n+1}+(n+1)\right\}-\left(\frac12a_n+n\right)\\
              &=\frac12a_{n+1}-\frac12a_n+1.
          \end{align*}
          したがって,
          \[
            a_{n+1}=-a_n+2.
          \]

          ここで $b_n=a_n-1$ とおくと,
          \begin{align*}
            b_{n+1}
              &=a_{n+1}-1\\
              &=-a_n+1\\
              &=-(a_n-1)\\
              &=-b_n.
          \end{align*}
          また,
          \[
            b_1=a_1-1=1
          \]
          であるから,
          \[
            b_n=(-1)^{n-1}.
          \]
          よって,
          \[
            a_n=1+(-1)^{n-1}.
          \]

      \end{enumerate}

    \subsubsection{純総和型 \, $S_{n+1}=pS_n+f(n) $}
      両辺に総和がきてもやることは一緒で,1つずらした総和と元の式を足し引きしてやることで,\,$a_n$や$a_{n+1}$の項を作れば良い.
      \begin{align*}
        &S_{n+1}=pS_n+f(n) \tag*{\(\cdots\) (1)}\\
        &S_{n}=pS_{n-1}+f(n-1) \tag*{\(\cdots\) (2)}\\
      \end{align*}
        $(1)-(2)$をして,
      \begin{align*}
        &a_{n+1}=pa_{n}+f(n)-f(n-1)
      \end{align*}
      あとはただの階差数列型の計算である.

      \noindent \textbf{【例題】}
      \begin{enumerate}[label=(\arabic*),parsep=3pt]
        \item $\displaystyle S_1=1,\ S_{n+1}=2S_n+n$
      \end{enumerate}

      \noindent\textbf{【方針】}
      \begin{enumerate}[label=(\arabic*),parsep=3pt]
        \item 本文の式は$n\ge2$でしか使えないことに注意し,\,$a_2$だけは$S_2=2S_1+f(1)$から別に求めましょう.
      \end{enumerate}

      \noindent\textbf{【解答】}
      \begin{enumerate}[label=(\arabic*),parsep=3pt]
        \item $\displaystyle a_1=1,\quad a_n=3\cdot2^{n-2}-1\ (n\ge2)$
      \end{enumerate}

      \noindent\textbf{【解法】}
      \begin{enumerate}[label=(\arabic*),parsep=3pt]
      \item $\displaystyle S_1=1,\ S_{n+1}=2S_n+n$
          $a_1=S_1=1$である.
          また,
          \begin{align*}
            &S_{n+1}=a_1+a_2+\cdots+a_n+a_{n+1}\\
            &S_n=a_1+a_2+\cdots+a_n
          \end{align*}
          であるから,
          \[
            S_{n+1}-S_n=a_{n+1}.
          \]
          したがって,
          \begin{align*}
            a_{n+1}
              &=S_{n+1}-S_n\\
              &=(2S_n+n)-S_n\\
              &=S_n+n.
          \end{align*}
          ここで,
          \[
            S_n=2S_{n-1}+(n-1)
          \]
          より,
          \begin{align*}
            a_{n+1}
              &=2S_{n-1}+(n-1)+n\\
              &=2S_{n-1}+2n-1.
          \end{align*}
          一方,
          \[
            a_n=S_n-S_{n-1}
          \]
          であるから,
          \[
            S_n=2S_{n-1}+(n-1)
          \]
          を用いて,
          \[
            a_n=S_{n-1}+n-1,
          \]
          したがって
          \[
            S_{n-1}=a_n-(n-1).
          \]
          これを上式に代入すると,
          \begin{align*}
            a_{n+1}
              &=2\{a_n-(n-1)\}+2n-1\\
              &=2a_n+1.
          \end{align*}

          なお,\,$n=1$のときは
          \[
            a_2=S_2-S_1=(2S_1+1)-S_1=2
          \]
          であり,上の漸化式
          \[
            a_{n+1}=2a_n+1
          \]
          は$n\ge1$で成り立つ.

          ここで $b_n=a_n+1$ とおくと,
          \begin{align*}
            b_{n+1}
              &=a_{n+1}+1\\
              &=2a_n+2\\
              &=2(a_n+1)\\
              &=2b_n.
          \end{align*}
          また,
          \[
            b_1=a_1+1=2
          \]
          であるから,
          \[
            b_n=2^n.
          \]
          よって,
          \[
            a_n=2^n-1.
          \]
      \end{enumerate}

